% =====================================================================
% Chapter 4: Verified population genetics
% Primary sources: P7 (doi:10.5281/zenodo.23653850), N4 (doi:10.5281/zenodo.23660071),
% H11/H12 (results/autoresearch/2026-10-25/), Lean: formal_cert/Autoresearch/PopGen*.lean,
% formal_cert/Programmes/PopGen/*.lean, formal_cert/Programmes/EvoGraph/*.lean.
% The Lean excerpts are copied verbatim, up to glyph rendering.
% =====================================================================
\chapter{Verified population genetics: fixation, selection and graphs}
\label{ch4:debut}\label{ch4:chapitre}
\chaptermark{Verified population genetics}
\sloppy

A mutation appears in a population of $N$ individuals. Will it disappear, or will it eventually be carried by
everyone? This is the question of \emph{fixation}. It is both very concrete (will a bacterium that has become
resistant to an antibiotic take over the colony?) and entirely accessible with the tools of a first course in
probability and analysis. All one needs is finite Markov chains, geometric sums, the binomial theorem, linear
recurrence relations and a little topology of $\R$.

This chapter follows the same thread as the previous ones. Each statement receives an undergraduate-level
\emph{human proof}, and next to it the statement \emph{as the Lean~4 kernel checked it}, with the file name and
line number. The two do not always coincide: sometimes the Lean proof takes a different route from the human
proof, and we point this out every time.

\begin{warnbox}[: status of this chapter]
\begin{itemize}[leftmargin=1.2em,itemsep=1pt]
\item All the results of this chapter are \textbf{classical}: Moran (1958) \cite{ch4:Mor}, Fisher, Wright,
Price (1970) \cite{ch4:Pri70}, Lieberman--Hauert--Nowak (2005) \cite{ch4:LHN}. No mathematical novelty is claimed. What is specific to the project
is a formalisation checked by the Lean kernel. We have not carried out a systematic search for earlier
formalisations.
\item The Lean files, the proofs and this text were produced by AI agents (Claude) under the author's
direction. A separate AI agent audited them: recompilation from source, axiom audit, planted controls. The
verdicts are recorded in \texttt{AUDIT.md} and \texttt{verification.md}.
\textbf{No human has audited these results.}
\item Nothing here depends on the openai/math theorems of Chapters~1 to~3.
\end{itemize}
\end{warnbox}

\paragraph{The tiers, recalled.} Here is what each label means in this chapter.
\begin{itemize}[itemsep=1pt]
\item \tierA{} means that the cited \emph{Lean statement} is checked by the kernel. Its axioms lie in
$\{\lean{propext},\lean{Classical.choice},\lean{Quot.sound}\}$, and the module has passed a kernel replay
(\texttt{leanchecker}) as well as a search for the soundness escapes \lean{skipKernelTC}, \lean{addDecl},
\lean{native_decide} and \lean{implemented_by}. Tier~A concerns the Lean statement, \emph{never} its
biological reading.
\item \tierL{} means a published source, read and located to the page.
\item \tierC{} means a human argument, including a proof in this book, that is not checked by
a machine.
\item \tierX{} means a computation (floating point, or exact rational arithmetic by a script) that illustrates and
checks, but never establishes anything.
\end{itemize}

\paragraph{Kernel replay coverage.} The original audits checked two things: the list of axioms, and
a recompilation by a separate agent. The N4 audit added a \texttt{leanchecker} replay of the modules
\lean{Programmes.PopGen.Absorption}, \lean{Programmes.PopGen.MoranModel} and the three modules
\lean{Programmes.EvoGraph.*}. Its negative control was a falsified theorem under \lean{debug.skipKernelTC},
and the replay did reject it. The other modules cited here had no replay: \lean{Autoresearch.PopGenFisher},
\lean{Autoresearch.PopGenWrightFisher}, \lean{Programmes.PopGen.WrightFisherFixation},
\lean{Programmes.PopGen.MoranFixation} and \lean{Programmes.PopGen.Price}. For these, the agent writing this
chapter ran \texttt{leanchecker} at the time of writing (2026-10-26), on the existing \texttt{.olean} files,
all newer than their sources, with the pinned Lean~4.34.1 toolchain; it also searched with \texttt{grep}
for the forbidden tokens. The result is given in §\ref{ch4:sec:bilan}. This replay is not a human audit.

\paragraph{The bridge to probability \tierC.} Throughout the chapter, $P$ denotes a stochastic matrix
on a finite state space. We define $P^n$ by the Chapman--Kolmogorov recursion $P^0=I$,
$P^{n+1}(i,j)=\sum_kP^n(i,k)P(k,j)$ (this is \lean{kPow} in Lean). If $(X_n)$ is the Markov chain with transition
matrix $P$ started at $i$, then by definition $\mathbb P_i(X_n=j)=P^n(i,j)$. If moreover the state $N$ is absorbing,
$\{X_n=N\}=\{T_N\le n\}$, where $T_N$ is the hitting time of $N$. The limit of $P^n(i,N)$ is therefore
$\mathbb P_i(T_N<\infty)$, the \emph{fixation probability}. This identification is standard, but it
\emph{is not} a Lean statement: no path space is constructed. The Lean theorems are about
limits of matrix powers, and that is the formulation we adopt, for instance ``the probability
of absorption at $N$ after $n$ steps tends to \dots''.

% ---------------------------------------------------------------------
\section{Hardy--Weinberg and one-locus selection}\label{ch4:sec:hw}\sectionmark{Hardy--Weinberg and selection}

\subsection{The model}
A gene (a \emph{locus}) has two alleles $A_1$ and $A_2$, with frequencies $p$ and $q=1-p$. Individuals are
diploid: each carries two copies, hence one of the three genotypes $A_1A_1$, $A_1A_2$ or $A_2A_2$.
Generations do not overlap. We further assume random mating: a zygote is formed from two
gametes drawn independently from the gamete pool.

\begin{proposition}[Hardy--Weinberg]\label{ch4:prop:hw}
Start from genotype frequencies $(x_{11},x_{12},x_{22})$ summing to~$1$, and set $p=x_{11}+x_{12}/2$. After one
generation of random mating, the frequencies are $(p^2,2pq,q^2)$, summing to~$1$, and the frequency of
allele $A_1$ is still $p$. These proportions are stationary from the first generation on.
\tierA{} (\lean{hwStep_sum}, \lean{hwStep_allele}, \lean{hwStep_stationary}, \texttt{PopGenFisher.lean},
ll.~129--148).
\end{proposition}

\begin{humanproof}
A gamete carries $A_1$ with probability $x_{11}+\frac12x_{12}=p$. Since the two gametes are independent, we obtain
$A_1A_1$ with probability $p^2$, $A_1A_2$ with probability $2pq$ (two possible orders) and $A_2A_2$ with
probability $q^2$. The sum is $(p+q)^2=1$. In the new generation, the frequency of $A_1$ is
$p^2+\frac12\cdot2pq=p(p+q)=p$: the map is therefore idempotent.
\end{humanproof}

\begin{leanbox}[: Hardy--Weinberg (\texttt{formal\_cert/Autoresearch/PopGenFisher.lean}, ll.~125--127 and 145)]
\begin{lstlisting}[language=Lean]
noncomputable def hwStep (x : ℝ × ℝ × ℝ) : ℝ × ℝ × ℝ :=
  let p := x.1 + x.2.1 / 2
  (p ^ 2, 2 * p * (1 - p), (1 - p) ^ 2)
theorem hwStep_stationary (x : ℝ × ℝ × ℝ) : hwStep (hwStep x) = hwStep x := by
\end{lstlisting}
The statement \lean{hwStep_stationary} holds for every triple, even an unnormalised one, because $p^2+p(1-p)=p$ is an identity.
On the other hand, \lean{hwStep} ignores $x_{22}$: it implicitly assumes a normalised input (remark by the
verifier, \texttt{H12/verification.md}).
\end{leanbox}

\subsection{Viability selection}
The genotypes $A_1A_1$, $A_1A_2$ and $A_2A_2$ survive to adulthood with \emph{fitnesses}
$w_{11},w_{12},w_{22}>0$. Before selection, the zygotes are in Hardy--Weinberg proportions. From this we
derive the \emph{mean fitness} and the new allele frequency:
\[
\bar w(p)=p^2w_{11}+2pq\,w_{12}+q^2w_{22},\qquad
p'=\frac{p^2w_{11}+pq\,w_{12}}{\bar w(p)}.
\]

\begin{leanbox}[: the model (\texttt{PopGenFisher.lean}, ll.~19--24)]
\begin{lstlisting}[language=Lean]
noncomputable def wbar (w11 w12 w22 p : ℝ) : ℝ :=
  p ^ 2 * w11 + 2 * p * (1 - p) * w12 + (1 - p) ^ 2 * w22

/-- Allele frequency after one generation of viability selection. -/
noncomputable def pnext (w11 w12 w22 p : ℝ) : ℝ :=
  (p ^ 2 * w11 + p * (1 - p) * w12) / wbar w11 w12 w22 p
\end{lstlisting}
\end{leanbox}

\begin{lemma}\label{ch4:lem:dp}
Let $w_{11},w_{12},w_{22}>0$ and $p\in[0,1]$. Then:
\begin{enumerate}[label=(\roman*),itemsep=1pt]
\item $\bar w(p)>0$ and $p'\in[0,1]$;
\item $\Delta p:=p'-p=\dfrac{pq\,s}{\bar w(p)}$, with $s=p(w_{11}-w_{12})+q(w_{12}-w_{22})$;
\item $\bar w'(p)=2s$, whence Wright's form $\Delta p=\dfrac{pq\,\bar w'(p)}{2\bar w(p)}$.
\end{enumerate}
\tierA{} (\lean{wbar_pos}, \lean{pnext_mem_Icc}, \lean{delta_p_eq}, \lean{hasDerivAt_wbar},
\lean{delta_p_eq_deriv}, \texttt{PopGenFisher.lean}, ll.~26--91).
\end{lemma}

\begin{humanproof}
(i) If $p>0$, the term $p^2w_{11}$ is $>0$ and the other two are $\ge0$. If $p=0$, $\bar w=w_{22}>0$. The
numerator of $p'$ is $\ge0$, and $\bar w-(p^2w_{11}+pqw_{12})=pqw_{12}+q^2w_{22}\ge0$, so $p'\le1$.
(ii) We bring everything over a common denominator, using $1-p=q$ and $1-2p=q-p$:
\[
p^2w_{11}+pqw_{12}-p\,\bar w
=p\bigl[pw_{11}(1-p)+qw_{12}(1-2p)-q^2w_{22}\bigr]
=pq\bigl[pw_{11}+(q-p)w_{12}-qw_{22}\bigr]=pq\,s ,
\]
then divide by $\bar w>0$.
(iii) Expanding, $\bar w(x)=w_{22}+2(w_{12}-w_{22})x+(w_{11}-2w_{12}+w_{22})x^2$, so
$\bar w'(p)=2(w_{12}-w_{22})+2(w_{11}-2w_{12}+w_{22})p=2s$.
\end{humanproof}

In Lean, the derivative in (iii) is not a hand-written polynomial. It is the real derivative \lean{deriv} of the
function $x\mapsto\bar w(x)$, obtained from an explicit \lean{HasDerivAt} statement (ll.~64--82).

\subsection{The exact increment identity}

\begin{theorem}[Monotonicity of mean fitness]\label{ch4:thm:wbar}
Let $w_{11},w_{12},w_{22}>0$ and $p\in[0,1]$. With the notation of Lemma~\ref{ch4:lem:dp}, we have the exact
identity
\begin{equation}\label{ch4:eq:increment}
\bar w(p')-\bar w(p)=\frac{pq\,s^2\,\bigl(\bar w(p)+p\,w_{11}+q\,w_{22}\bigr)}{\bar w(p)^2},
\end{equation}
and consequently $\bar w(p')\ge\bar w(p)$: mean fitness never decreases from one generation to the
next. \tierA{} (\lean{wbar_increment} and \lean{fisher_mean_fitness_nondecreasing},
\texttt{PopGenFisher.lean}, ll.~95--119).
\end{theorem}

\begin{humanproof}
Set $W=\bar w(p)$, $a=w_{11}-2w_{12}+w_{22}$ and $\Delta=p'-p=pqs/W$ (Lemma~\ref{ch4:lem:dp}).

\emph{Step 1: Taylor is exact.} $\bar w$ is a polynomial of degree~$2$ in $x$, with derivative $2s$ at $p$ and
second derivative $2a$. Taylor's formula to order~$2$ is therefore exact:
\[
\bar w(p+\Delta)=W+2s\,\Delta+a\,\Delta^2 .
\]

\emph{Step 2: substitute $\Delta$.}
\[
\bar w(p')-W=2s\cdot\frac{pqs}{W}+a\,\frac{p^2q^2s^2}{W^2}=\frac{pq\,s^2}{W^2}\,\bigl(2W+a\,pq\bigr).
\]

\emph{Step 3: the key identity.} Expanding, one checks that
\[
W+a\,pq=p^2w_{11}+2pqw_{12}+q^2w_{22}+pq(w_{11}-2w_{12}+w_{22})=pw_{11}(p+q)+qw_{22}(p+q)=pw_{11}+qw_{22}.
\]
Thus $2W+a\,pq=W+pw_{11}+qw_{22}$, which gives exactly~\eqref{ch4:eq:increment}.

\emph{Step 4: the sign.} The factors $p$, $q$, $s^2$, $W^2$ are $\ge0$ (and $W>0$), and
$W+pw_{11}+qw_{22}>0$. The quotient is therefore $\ge0$.
\end{humanproof}

\begin{leanbox}[: the increment (\texttt{PopGenFisher.lean}, ll.~95--99 and 110--112)]
\begin{lstlisting}[language=Lean]
theorem wbar_increment {w11 w12 w22 p : ℝ} (h11 : 0 < w11) (h12 : 0 < w12) (h22 : 0 < w22)
    (hp0 : 0 ≤ p) (hp1 : p ≤ 1) :
    wbar w11 w12 w22 (pnext w11 w12 w22 p) - wbar w11 w12 w22 p
      = p * (1 - p) * (p * (w11 - w12) + (1 - p) * (w12 - w22)) ^ 2
          * (wbar w11 w12 w22 p + p * w11 + (1 - p) * w22) / (wbar w11 w12 w22 p) ^ 2 := by
theorem fisher_mean_fitness_nondecreasing {w11 w12 w22 p : ℝ}
    (h11 : 0 < w11) (h12 : 0 < w12) (h22 : 0 < w22) (hp0 : 0 ≤ p) (hp1 : p ≤ 1) :
    wbar w11 w12 w22 p ≤ wbar w11 w12 w22 (pnext w11 w12 w22 p) := by
\end{lstlisting}
\emph{Side by side.} The Lean proof does not go through Taylor. It clears denominators (\lean{field_simp}),
then concludes with \lean{ring}. The docstring states that the identity was first found with \texttt{sympy}, then
checked by the kernel. The human proof above explains \emph{why} the factor
$\bar w+pw_{11}+qw_{22}$ appears: it is the identity of Step~3.
\end{leanbox}

\begin{remark}[Hypotheses \tierC]\label{ch4:rem:positivite}
\emph{Strict} positivity of the fitnesses is sufficient, but not necessary: by~\eqref{ch4:eq:increment},
$w\ge0$ with $\bar w(p)>0$ suffices, since then $\bar w+pw_{11}+qw_{22}\ge\bar w>0$. This is a correction by the
independent verifier (\texttt{H12/verification.md}). The Lean theorem is correct, but narrower than it
could be. \emph{Non-negativity}, however, cannot be dropped. The control
\lean{fisher_fails_without_positivity} takes $w=(-1,1,0)$ and $p=\frac12$: then $\bar w=\frac14>0$, but $p'=0$ and
$\bar w(p')=0<\frac14$ \tierA{} (ll.~164--167).
\end{remark}

\begin{remark}[The mean goes up, but $p$ may move away from equilibrium \tierA]
Under underdominance, $w=(1,\frac12,1)$, there is an interior equilibrium $p=\frac12$. Yet, starting from
$p=\frac25$, we get $p'=\frac7{19}<\frac25$: the frequency moves \emph{away} from the equilibrium, and $\bar w$
increases all the same ($\frac{19}{25}<\frac{277}{361}$). This is \lean{underdominance_moves_away} (ll.~155--158),
and the numerical values were recomputed in exact rational arithmetic \tierX. The increase of $\bar w$ therefore says
nothing about convergence to a polymorphism.
\end{remark}

\begin{warnbox}[: this is not Fisher's ``fundamental theorem'']
Theorem~\ref{ch4:thm:wbar} covers a single locus, two alleles, discrete generations, viability
selection, constant fitnesses and random mating. It is the monotonicity of mean
fitness \emph{for this model}. Fisher's fundamental theorem, in the sense of Price's (1972) reading
in terms of the ``partial change'' \cite{ch4:Pri72} and additive genetic variance, is a different, finer statement, which
\emph{is not formalised} here \tierC. The original Lean docstring spoke of ``Fisher's fundamental theorem''. The
verifier required it to be renamed, and the module header now says ``one-locus, two-allele mean-fitness
monotonicity''.
\end{warnbox}

% ---------------------------------------------------------------------
\section[Neutral Wright--Fisher]{Neutral Wright--Fisher: martingale, harmonic functions, absorption}\label{ch4:sec:wf}

\subsection{The chain}
We now turn to \emph{chance} alone. Consider a population of $N$ genes, of which $i$ are of mutant
type. The next generation is obtained by drawing $N$ genes \emph{with replacement} from the present generation.
Without selection, the number $j$ of mutants in the next generation follows the binomial distribution
$\mathcal B(N,i/N)$:
\[
P(i,j)=\binom Nj\Bigl(\frac iN\Bigr)^j\Bigl(1-\frac iN\Bigr)^{N-j},\qquad 0\le i,j\le N .
\]

\begin{leanbox}[: the transition kernel (\texttt{formal\_cert/Autoresearch/PopGenWrightFisher.lean}, ll.~25--26)]
\begin{lstlisting}[language=Lean]
noncomputable def wfP (N : ℕ) (i j : Fin (N + 1)) : ℝ :=
  (N.choose (j : ℕ) : ℝ) * ((i : ℕ) / N : ℝ) ^ (j : ℕ) * (1 - ((i : ℕ) / N : ℝ)) ^ (N - (j : ℕ))
\end{lstlisting}
The positive control \lean{ctl_positive_wfP_value} establishes \lean{wfP 2 1 1 = 1 / 2}: the division is indeed
real division, not the integer division of $\N$.
\end{leanbox}

\begin{proposition}\label{ch4:prop:wfbase}
Let $N\ge1$.
\begin{enumerate}[label=(\roman*),itemsep=1pt]
\item $P$ is stochastic, and the states $0$ and $N$ are absorbing.
\item \emph{(Martingale)} For every $i$, $\sum_jP(i,j)\,j=i$. In other words, $h(j)=j/N$ is \emph{harmonic}:
$Ph=h$.
\item From every state $0\le i<N$, we have $P(i,0)=(1-i/N)^N>0$.
\end{enumerate}
\tierA{} (\lean{wfP_nonneg}, \lean{wfP_sum_one}, \lean{wfP_zero}, \lean{wfP_last}, \lean{wfP_mean},
\lean{frac_harmonic}, \lean{wfP_interior_zero_pos}, \texttt{PopGenWrightFisher.lean}, ll.~34--118).
\end{proposition}

\begin{humanproof}
(i) By the binomial theorem, $\sum_j\binom Nj x^j(1-x)^{N-j}=(x+1-x)^N=1$. At $i=0$ we have $x=0$, and with the
convention $0^0=1$ the row is $\delta_0$; likewise at $i=N$.
(ii) We use $j\binom Nj=N\binom{N-1}{j-1}$ for $j\ge1$. With $x=i/N$ and $y=1-x$,
\[
\sum_{j=0}^N\binom Nj x^jy^{N-j}\,j=N\sum_{j=1}^{N}\binom{N-1}{j-1}x^jy^{N-j}=Nx\,(x+y)^{N-1}=Nx=i .
\]
(iii) This is the term $j=0$, which equals $(1-i/N)^N$, strictly positive as soon as $i<N$.
\end{humanproof}

Identity (ii) says that the number $X_n$ of mutants is a \emph{martingale}: its conditional expectation at the
next step is its present value \tierC{} (probabilistic reading). By induction on $n$, we obtain
$\sum_jP^n(i,j)\,j=i$ for every $n$. This computation requires no measure theory.

\begin{uniremark}\tierL{} \emph{Optional stopping.} In the path-space language, $(X_n)$ is a martingale with
values in $\{0,\dots,N\}$, hence bounded, and the absorption time $T=\min(T_0,T_N)$ is almost surely finite (this
follows, for instance, from Lemma~\ref{ch4:lem:geom} below). The optional stopping theorem for bounded martingales
then gives $\mathbb E_i[X_T]=X_0=i$, that is $N\,\mathbb P_i(X_T=N)=i$, so the fixation probability is $i/N$.
This is the standard martingale proof \cite{app:williams,app:norris}; it is \emph{not} the route taken in Lean,
which never constructs the path space.
\end{uniremark}

\subsection{Harmonic functions and the maximum principle}
\begin{definition}
Let $P$ be a stochastic matrix on $\{0,\dots,N\}$ whose states $0$ and $N$ are absorbing. A function
$g:\{0,\dots,N\}\to\R$ is \emph{harmonic in the interior} if $\sum_jP(i,j)g(j)=g(i)$ for every $0<i<N$.
\end{definition}

The point of this notion is the following. The fixation probability $u(i)=\mathbb P_i(T_N<\infty)$ satisfies
\emph{formally} the ``first-step'' equation $u(i)=\sum_jP(i,j)u(j)$, with $u(0)=0$ and $u(N)=1$ \tierC.
If one knows that only one harmonic function has these boundary values, one therefore knows $u$.

\begin{uniremark}\tierL{} \emph{Hitting probabilities as harmonic functions.} For a Markov chain on a countable
state space and a target set $A$, the vector of hitting probabilities $u(i)=\mathbb P_i(\text{hit }A)$ is the
\emph{minimal} non-negative solution of the system $u=1$ on $A$, $u(i)=\sum_jP(i,j)u(j)$ off $A$; on a
finite state space from which $\{0,N\}$ is reached almost surely, the bounded solution with prescribed boundary
values is unique, which is exactly what makes the first-step equation determine $u$ \cite{app:norris} (standard).
\end{uniremark}

\begin{theorem}[Maximum principle, uniqueness]\label{ch4:thm:maxprinc}
Let $N\ge1$, and let $d$ be a function harmonic in the interior for the Wright--Fisher kernel, with $d(0)=d(N)=0$.
Then $d\le0$, and by applying the same reasoning to $-d$, $d=0$. Consequently, $i\mapsto i/N$ is
\emph{the unique} function harmonic in the interior that equals $0$ at $0$ and $1$ at $N$. \tierA{} (\lean{max_principle},
\lean{harmonic_unique}, \lean{wright_fisher_fixation_harmonic}, \texttt{PopGenWrightFisher.lean}, ll.~122--187).
\end{theorem}

\begin{humanproof}
Let $m$ be a point where $d$ attains its maximum; it exists because the set is finite. If $m\in\{0,N\}$, then
$\max d=0$. Otherwise $m$ is interior, and
\[
0=d(m)-\sum_jP(m,j)d(j)=\sum_jP(m,j)\bigl(d(m)-d(j)\bigr),
\]
where we used $\sum_jP(m,j)=1$. Each term is $\ge0$, since $P\ge0$ and $d(m)$ is maximal; they
are therefore all zero. But $P(m,0)>0$ (Proposition~\ref{ch4:prop:wfbase}(iii)), so $d(m)=d(0)=0$. In both
cases, $\max d=0$. For uniqueness, if $g$ is a solution, then $d=g-h$ with $h(j)=j/N$ is harmonic in the interior and
vanishes on the boundary.
\end{humanproof}

\begin{leanbox}[: the maximum principle (\texttt{PopGenWrightFisher.lean}, ll.~122--126)]
\begin{lstlisting}[language=Lean]
theorem max_principle {N : ℕ} (hN : 0 < N) (d : Fin (N + 1) → ℝ) (h0 : d 0 = 0)
    (hN' : d (Fin.last N) = 0)
    (hharm : ∀ i : Fin (N + 1), 0 < (i : ℕ) → (i : ℕ) < N → d i = ∑ j, wfP N i j * d j) :
    ∀ i, d i ≤ 0 := by
\end{lstlisting}
The Lean proof follows the human proof exactly. \lean{Finset.exists_max_image} supplies the maximum point, then
\lean{Finset.sum_eq_zero_iff_of_nonneg} kills each term, and \lean{wfP_interior_zero_pos} concludes.
\end{leanbox}

\begin{warnbox}[: what H11 did not prove]
Module H11 \cite{ch4:H1112} stops there. It establishes that $i/N$ is \emph{the} harmonic solution, but the step ``the fixation
probability is harmonic'' is not formalised in it. The verifier put it bluntly: ``this file does NOT
kernel-check `fixation probability $= i/N$'\,'' (\texttt{H11/NOTE.md}). It suggested closing the gap
\emph{without} measure theory, by proving $P^n(i,N)\to i/N$. This is what programme P7 did, and it is
the subject of the next subsection.
\end{warnbox}

\subsection{Absorption: \texorpdfstring{$P^n(i,N)\to i/N$}{P\^{}n(i,N) -> i/N}}
We write $I_n(i)=\sum_{0<k<N}P^n(i,k)$ for the \emph{interior mass} after $n$ steps, that is, the probability of not
yet having been absorbed (\lean{interiorMass}).

\begin{theorem}[Neutral fixation]\label{ch4:thm:wflim}
For every $N\ge1$ and every $0\le i\le N$,
\[
P^n(i,N)\xrightarrow[n\to\infty]{}\frac iN,\qquad P^n(i,0)\xrightarrow[n\to\infty]{}1-\frac iN,
\]
and the sequence $\bigl(P^n(i,N)\bigr)_n$ is nondecreasing. \tierA{} (\lean{wright_fisher_fixation_limit},
\lean{wright_fisher_fixation_mono}, \texttt{Programmes/PopGen/}\allowbreak\texttt{WrightFisherFixation.lean}, ll.~49--58). The general theorem \lean{fixation_limit} from which it follows reappears under the name \lean{KAL}.
\end{theorem}

We first give an undergraduate-level proof \tierC, which relies on the \emph{geometric decay}
of the interior mass. The Lean proof takes a different route, presented afterwards.

\begin{lemma}[Geometric decay \tierC]\label{ch4:lem:geom}
Set $\delta=\min_{0<k<N}P(k,0)$. Then $\delta\ge N^{-N}>0$ and, for every $n$, $I_{n+1}(i)\le(1-\delta)\,I_n(i)$.
Consequently $I_n(i)\le(1-\delta)^n$.
\end{lemma}

\begin{humanproof}
For $0<k<N$, we have $P(k,0)=(1-k/N)^N\ge(1/N)^N$. Next,
\[
I_{n+1}(i)=\sum_{0<j<N}\sum_kP^n(i,k)P(k,j)=\sum_kP^n(i,k)\sum_{0<j<N}P(k,j).
\]
If $k\in\{0,N\}$, the inner sum is zero, because the state is absorbing. If $k$ is interior, it is at most
$1-P(k,0)\le1-\delta$. Hence $I_{n+1}(i)\le(1-\delta)\sum_{0<k<N}P^n(i,k)=(1-\delta)I_n(i)$, and we conclude by
induction, since $I_0\le1$.
\end{humanproof}

\begin{humanproof}[Human proof of Theorem~\ref{ch4:thm:wflim}]
Set $h(j)=j/N$. By Proposition~\ref{ch4:prop:wfbase}, $Ph=h$ at every state, hence $P^nh=h$, which reads
\[
\frac iN=\sum_jP^n(i,j)\,\frac jN=P^n(i,N)+\sum_{0<j<N}P^n(i,j)\,\frac jN .
\]
The last term lies between $0$ and $I_n(i)$, which tends to $0$ by Lemma~\ref{ch4:lem:geom}. We deduce
$P^n(i,N)\to i/N$ by squeezing. Next, $P^n(i,0)+P^n(i,N)+I_n(i)=1$ (the rows of $P^n$ sum to~$1$),
whence $P^n(i,0)\to1-i/N$. Finally, $P^{n+1}(i,N)=\sum_kP^n(i,k)P(k,N)\ge P^n(i,N)P(N,N)=P^n(i,N)$.
\end{humanproof}

\begin{uniremark}\tierL{} \emph{The fundamental matrix.} Order the states so that the interior states
$\{1,\dots,N-1\}$ come first and write $P=\begin{pmatrix}Q&R\\0&I\end{pmatrix}$, where $Q$ is the interior-to-interior
block and $R$ the interior-to-boundary block. Lemma~\ref{ch4:lem:geom} says that $Q^n\to0$ geometrically, so
the series $\sum_{n\ge0}Q^n$ converges and equals the \emph{fundamental matrix} $N_Q=(I-Q)^{-1}$; its entry
$N_Q(i,k)$ is the expected number of visits to $k$ starting from $i$. The absorption probabilities are the entries
of $B=N_QR$, and the vector of expected absorption times is $N_Q\mathbf 1$ \cite{app:norris} (standard). This is
literature only: as recorded in §\ref{ch4:sec:bilan}, absorption times are not formalised in this chapter.
\end{uniremark}

\begin{leanbox}[: the general theorem and the Wright--Fisher case]
\texttt{formal\_cert/Programmes/PopGen/Absorption.lean}, ll.~32--34, 45--53 and 302--308:
\begin{lstlisting}[language=Lean]
noncomputable def kPow (P : Fin (N + 1) → Fin (N + 1) → ℝ) : ℕ → Fin (N + 1) → Fin (N + 1) → ℝ
  | 0 => fun i j => if i = j then 1 else 0
  | n + 1 => fun i j => ∑ k, kPow P n i k * P k j
structure AbsorbingChain (P : Fin (N + 1) → Fin (N + 1) → ℝ) : Prop where
  nonneg : ∀ i j, 0 ≤ P i j
  rowSum : ∀ i, ∑ j, P i j = 1
  absorb0 : ∀ j, P 0 j = if j = 0 then 1 else 0
  absorbN : ∀ j, P (Fin.last N) j = if j = Fin.last N then 1 else 0
def InteriorHarmonic (P : Fin (N + 1) → Fin (N + 1) → ℝ) (g : Fin (N + 1) → ℝ) : Prop :=
  ∀ i : Fin (N + 1), 0 < (i : ℕ) → (i : ℕ) < N → ∑ j, P i j * g j = g i
theorem fixation_limit {P : Fin (N + 1) → Fin (N + 1) → ℝ} (hP : AbsorbingChain P) (hN : 0 < N)
    {h : Fin (N + 1) → ℝ} (hh : InteriorHarmonic P h) (h0 : h 0 = 0) (h1 : h (Fin.last N) = 1)
    (hvar : ∀ k : Fin (N + 1), 0 < (k : ℕ) → (k : ℕ) < N → 0 < ∑ j, P k j * (h j - h k) ^ 2)
    (i : Fin (N + 1)) :
    Tendsto (fun n => kPow P n i (Fin.last N)) atTop (𝓝 (h i)) ∧
    Tendsto (fun n => kPow P n i 0) atTop (𝓝 (1 - h i)) := by
\end{lstlisting}
\texttt{WrightFisherFixation.lean}, ll.~49--51:
\begin{lstlisting}[language=Lean]
theorem wright_fisher_fixation_limit {N : ℕ} (hN : 0 < N) (i : Fin (N + 1)) :
    Tendsto (fun n => kPow (wfP N) n i (Fin.last N)) atTop (𝓝 (((i : ℕ) : ℝ) / N)) ∧
    Tendsto (fun n => kPow (wfP N) n i 0) atTop (𝓝 (1 - ((i : ℕ) : ℝ) / N)) := by
\end{lstlisting}
\emph{Side by side: two different routes.} The Lean proof does not use Lemma~\ref{ch4:lem:geom}. It
uses a \emph{second-moment} argument, valid for every absorbing chain provided that a harmonic
function $h$ has a one-step variance $v(k)=\sum_jP(k,j)(h(j)-h(k))^2$ that is strictly positive at every interior
state (hypothesis \lean{hvar}). Here are the steps.
\begin{enumerate}[itemsep=1pt]
\item We have the identity $\sum_jP(k,j)h(j)^2=h(k)^2+v(k)$ (\lean{second_moment}, l.~145).
\item With $\varepsilon=\min v>0$, we obtain $P^{m+1}(h^2)\ge P^m(h^2)+\varepsilon I_m$. Summing,
$\varepsilon\sum_{m<n}I_m(i)\le\sum_jh(j)^2$.
\item The series $\sum I_m$ therefore has non-negative terms and bounded partial sums: it converges, and
$I_m\to0$ (\lean{interiorMass_tendsto_zero}, l.~156).
\item We conclude by the squeeze above (\lean{absorption_tendsto}, \lean{mass_split}).
\end{enumerate}
This route does not give the geometric rate. On the other hand, it applies as is to the Moran model
(§\ref{ch4:sec:moran}) and to graphs (§\ref{ch4:sec:graphes}), where a direct jump to $0$ is impossible. For
Wright--Fisher, the variance is positive thanks to the jump to~$0$ alone: $v(k)\ge P(k,0)(k/N)^2>0$
(\lean{wf_var_pos}, l.~34).
\end{leanbox}

\begin{remark}[Controls \tierA]\label{ch4:rem:wfctl}
\emph{Positive control.} We have $P^2(1,2)=\frac38$ for $N=2$ (\lean{ctl_pos_wf_two_steps}, l.~62). The limit
$\frac12$ is therefore not yet reached after two generations.
\emph{Negative control.} The variance hypothesis cannot be dropped. The identity matrix on $\{0,1,2\}$
is stochastic and absorbing, and $j/2$ is harmonic for it; yet $P^n(1,2)=0$ for every $n$
(\lean{ctl_neg_identity_no_absorption}, \texttt{Absorption.lean}, l.~363).
\emph{Uniqueness.} If $g$ is another harmonic function with the same boundary values, the same argument gives
$P^n(i,N)\to g(i)$. By uniqueness of limits, $g=h$ (\lean{harmonic_unique_of_var}, l.~320). This replaces the
maximum principle by the uniqueness of limits.
\end{remark}

\begin{remark}[The Genetic repository]\label{ch4:rem:genetic}
The companion repository \texttt{SocrateAI-Scientific-Genetic} \cite{ch4:Genetic} contains an earlier version of these results, over
$\mathbb Q$. It is in \texttt{formal/ANSE/}\allowbreak\texttt{AlixPopulationGenetics.lean}: the maximum principle
\lean{ilene_harmonic_le_zero} (l.~82), uniqueness \lean{xav_wf_harmonic_solution_unique} (l.~199) and the bound
without passing to the limit \lean{xav_wf_fixation_by_time_le} (l.~268), which asserts $P^t(i,N)\le i/N$ for every $t$.
The latter follows in one line from the martingale:
$P^t(i,N)\cdot1\le\sum_jP^t(i,j)\,j/N=i/N$. The proof of this one-line bound is a human proof \tierC. The file
was not recompiled for this book, and we assign it no tier. The statements of this chapter
come from the AutoevolveAI modules H11 and P7. Between the two repositories, the link is only textual: no Lean
term connects the two namespaces \tierC.
\end{remark}

% ---------------------------------------------------------------------
\section{The Moran process with selection}\label{ch4:sec:moran}\sectionmark{Moran with selection}

\subsection{The model}
A population has $N$ individuals, of which $i$ are mutants with fitness $r>0$; residents have
fitness $1$. At each time step, two independent draws take place.
\begin{itemize}[itemsep=1pt]
\item One individual is chosen to reproduce, with probability proportional to its fitness. A
mutant is therefore chosen with probability $\frac{ri}{ri+N-i}$.
\item One individual is chosen \emph{uniformly} among the $N$ to die, and the offspring takes its place.
\end{itemize}
The number of mutants goes up by $1$ with probability $p_i$, goes down by $1$ with probability $q_i$ and stays
unchanged otherwise:
\[
p_i=\frac{ri}{ri+N-i}\cdot\frac{N-i}{N},\qquad q_i=\frac{N-i}{ri+N-i}\cdot\frac iN .
\]
These are the rates of Keller and Uğurlu, up to a one-line rewriting \tierL{}~\cite{ch4:KU} (pp.~2--4) \tierC.

\begin{leanbox}[: the Moran model (\texttt{formal\_cert/Programmes/PopGen/MoranModel.lean}, ll.~30--32, 74--83, 152--155)]
\begin{lstlisting}[language=Lean]
noncomputable def bdP (N : ℕ) (p q : ℕ → ℝ) (i j : Fin (N + 1)) : ℝ :=
  (if (j : ℕ) = (i : ℕ) + 1 then p i else 0) + (if (j : ℕ) + 1 = (i : ℕ) then q i else 0) +
    (if (j : ℕ) = (i : ℕ) then 1 - p i - q i else 0)
noncomputable def moranUp (N : ℕ) (r : ℝ) (i : ℕ) : ℝ :=
  r * i / (r * i + ((N : ℝ) - i)) * (((N : ℝ) - i) / N)
noncomputable def moranDown (N : ℕ) (r : ℝ) (i : ℕ) : ℝ :=
  ((N : ℝ) - i) / (r * i + ((N : ℝ) - i)) * ((i : ℝ) / N)
noncomputable def moranP (N : ℕ) (r : ℝ) : Fin (N + 1) → Fin (N + 1) → ℝ :=
  bdP N (moranUp N r) (moranDown N r)
noncomputable def geomS (r : ℝ) (n : ℕ) : ℝ := ∑ k ∈ range n, r⁻¹ ^ k
noncomputable def moranFix (N : ℕ) (r : ℝ) (i : ℕ) : ℝ := geomS r i / geomS r N
\end{lstlisting}
The independent auditor checked symbolically that \lean{moranUp} and \lean{moranDown} coincide with the
Keller--Uğurlu rates (PG-C-0004). This is an AI judgement, not a human audit.
\end{leanbox}

\begin{lemma}\label{ch4:lem:moranbase}
Let $N\ge1$ and $r>0$.
\begin{enumerate}[label=(\roman*),itemsep=1pt]
\item $p_i=r\,q_i$ for every $i$: the ratio of the rates is the fitness.
\item $1-p_i-q_i=\dfrac{ri^2+(N-i)^2}{(ri+N-i)N}\ge0$, so the matrix is stochastic. Moreover
$p_0=q_0=p_N=q_N=0$, so $0$ and $N$ are absorbing.
\item For $0<i<N$, we have $p_i>0$ and $q_i>0$.
\end{enumerate}
\tierA{} (\lean{moranUp_eq}, \lean{moran_stay_eq}, \lean{moran_stay_nonneg}, \lean{moranUp_pos},
\lean{moranDown_pos}, \lean{moran_absorbing}, \texttt{MoranModel.lean}, ll.~85--149, and
\texttt{MoranFixation.lean}, l.~25).
\end{lemma}

\begin{humanproof}
(i) We have $p_i/q_i=\frac{ri(N-i)}{(N-i)i}=r$ if $0<i<N$, and both rates vanish at the boundary.
(ii) We compute
$(ri+N-i)N-ri(N-i)-(N-i)i=riN-ri(N-i)+(N-i)(N-i)=ri^2+(N-i)^2$.
(iii) All factors are strictly positive when $0<i<N$.
\end{humanproof}

\subsection{Gambler's ruin}

\begin{theorem}[Fixation in the Moran model]\label{ch4:thm:moran}
Let $N\ge1$, $r>0$, and $S_n=\sum_{k<n}r^{-k}$. For every $0\le i\le N$,
\[
P^n(i,N)\xrightarrow[n\to\infty]{}\frac{S_i}{S_N}=
\begin{cases}\dfrac{1-r^{-i}}{1-r^{-N}}&\text{if }r\ne1,\\[2mm] \dfrac iN&\text{if }r=1,\end{cases}
\qquad P^n(i,0)\longrightarrow1-\frac{S_i}{S_N}.
\]
Moreover, $S_i/S_N$ is the unique function harmonic in the interior that equals $0$ at $0$ and $1$ at $N$, and
$\frac{1-r^{-i}}{1-r^{-N}}\to\frac iN$ as $r\to1$, $r\ne1$. \tierA{} (\lean{moran_fixation_limit},
\lean{moran_fixation_limit_closed}, \lean{moran_fixation_limit_neutral}, \lean{moran_fixation_unique},
\texttt{MoranFixation.lean}, ll.~72--97; \lean{moranFix_closed}, \lean{moran_closed_tendsto_neutral},
\texttt{MoranModel.lean}, ll.~203 and 231).
\end{theorem}

\begin{humanproof}
\emph{Step 1: the recurrence.} Let $h$ be harmonic in the interior, with $h(0)=0$ and $h(N)=1$. For $0<i<N$,
the first-step equation reads
\[
p_ih(i+1)+q_ih(i-1)+(1-p_i-q_i)h(i)=h(i)\iff p_i\bigl(h(i+1)-h(i)\bigr)=q_i\bigl(h(i)-h(i-1)\bigr).
\]
Since $p_i=rq_i$ and $q_i>0$, the increments $\Delta_i=h(i+1)-h(i)$ satisfy $\Delta_i=r^{-1}\Delta_{i-1}$.
This is a geometric sequence: $\Delta_i=r^{-i}\Delta_0$. Summing, $h(i)=\Delta_0\,S_i$. The condition $h(N)=1$
then forces $\Delta_0=1/S_N$, with $S_N>0$. There is therefore \emph{at most} one solution, and conversely
$h(i)=S_i/S_N$ is one, since the equivalences can be reversed. This is the classical ``gambler's
ruin'' argument.

\emph{Step 2: closed form.} If $r\ne1$, set $\rho=1/r\ne1$. Then $S_n=\frac{1-\rho^n}{1-\rho}$, whence
$\frac{S_i}{S_N}=\frac{1-r^{-i}}{1-r^{-N}}$. If $r=1$, $S_n=n$. Each $S_n(r)$ is a finite sum of continuous
functions of $r$ and $S_N(1)=N\ne0$, so the quotient is continuous at $r=1$. This is the neutral limit.

\emph{Step 3: absorption.} We repeat the proof of Theorem~\ref{ch4:thm:wflim} with $h=S_\cdot/S_N$:
$P^nh=h$, then $|h(i)-P^n(i,N)|\le I_n(i)$, because $0\le h\le1$. It remains to show that $I_n(i)\to0$.
Lemma~\ref{ch4:lem:geom} no longer applies as is: from $i$, one cannot jump to $0$ in one step if
$i\ge2$. We therefore argue in blocks. From every interior state $k$, the path $k\to k-1\to\dots\to0$ has
probability $q_kq_{k-1}\cdots q_1>0$ and lasts only $k\le N-1$ steps. Once at $0$, the chain stays there. Set
$m=N-1$ and $\eta=\min_{0<k<N}q_k\cdots q_1>0$. Then $P^m(k,0)\ge\eta$ for every interior $k$, and the same
computation as in Lemma~\ref{ch4:lem:geom}, applied to the matrix $P^m$, gives $I_{n+m}(i)\le(1-\eta)I_n(i)$.
Since $I_n$ is nonincreasing (absorbed mass stays absorbed), $I_n(i)\to0$ \tierC.
\end{humanproof}

\begin{leanbox}[: the theorem and its closed form (\texttt{MoranFixation.lean}, ll.~72--74; \texttt{MoranModel.lean}, ll.~173--175, 203--204)]
\begin{lstlisting}[language=Lean]
theorem moran_fixation_limit {N : ℕ} {r : ℝ} (hN : 0 < N) (hr : 0 < r) (i : Fin (N + 1)) :
    Tendsto (fun n => kPow (moranP N r) n i (Fin.last N)) atTop (𝓝 (moranFix N r i)) ∧
    Tendsto (fun n => kPow (moranP N r) n i 0) atTop (𝓝 (1 - moranFix N r i)) :=
theorem moranFix_harmonic_nat {N : ℕ} {r : ℝ} (hr : 0 < r) {i : ℕ} (h1 : 0 < i) :
    moranUp N r i * moranFix N r (i + 1) + moranDown N r i * moranFix N r (i - 1)
      + (1 - moranUp N r i - moranDown N r i) * moranFix N r i = moranFix N r i := by
theorem moranFix_closed {N : ℕ} {r : ℝ} (hr1 : r ≠ 1) (i : ℕ) :
    moranFix N r i = (1 - r⁻¹ ^ i) / (1 - r⁻¹ ^ N) := by
\end{lstlisting}
\emph{Side by side.} Lean does not \emph{solve} the recurrence. It \emph{checks} that the guessed solution
$S_i/S_N$ is harmonic (\lean{moranFix_harmonic_nat}, with a single \lean{linear_combination}). It then applies the
general theorem \lean{fixation_limit} of §\ref{ch4:sec:wf}, with the variance
$v(i)\ge p_i\,(r^{-i}/S_N)^2>0$ (\lean{moran_var_pos_nat}, l.~189). The block argument of Step~3 is not
formalised; the second moment replaces it. Uniqueness comes from the uniqueness of limits
(\lean{moran_fixation_unique}), not from Step~1.
\end{leanbox}

\begin{remark}[Controls \tierA]\label{ch4:rem:moranctl}
The following controls are all checked by the kernel.
\begin{itemize}[itemsep=1pt]
\item \emph{Positive controls.} For $N=2$, $r=2$, we have $p_1=\frac13$, $q_1=\frac16$ and $h(1)=\frac23$
(\lean{ctl_pos_moran_rates}, \lean{ctl_pos_moranFix_N2}). For $N=3$, $r=2$, we have $h=(0,\frac47,\frac67,1)$, and
the first-step equation holds (\lean{ctl_pos_moranFix_N3}, \texttt{MoranModel.lean}, ll.~240--253).
\item \emph{Negative controls.} The ``inverted-exponent'' formula $(1-r^i)/(1-r^N)$ gives $\frac17\ne\frac47$
and violates the first-step equation. The neutral value $i/N$ also violates it under selection (ll.~259--276). At the level
of the chain, $P^n(1,3)$ for $N=3$, $r=2$ tends neither to $\frac13$ nor to $\frac17$
(\lean{ctl_neg_wrong_limits}, \texttt{MoranFixation.lean}, l.~102).
\item \emph{Planted controls.} In a separate file that nothing imports, a ``theorem'' proved by \lean{sorry}
and another derived from a planted axiom assert the false limits $\frac17$ and $\frac13$. This file
(\texttt{PopGenProgControls.lean}) compiles with return code~$0$, but \texttt{\#print axioms} shows \lean{sorryAx}, resp.\ the planted axiom
\lean{ctl_cheat}. They are therefore detected (PG-A-0012).
\end{itemize}
\end{remark}

\begin{numbox}[: CVODE and the ``fail-closed'' certificate]
\tierX{} Programme P7 also integrated the uniformised Kolmogorov equation $\dot\pi=\pi(P-I)$ with CVODE (BDF, exact
Jacobian), in a new crate \texttt{qf-popgen} of rusty-SUNDIALS. The crate wraps itself in a certificate (``fail-closed'': it refuses by default) that rejects
any doubtful result: valid kernel, mass conservation, componentwise positivity,
decay of the interior mass and agreement with the target. Outcome: $106$ accepted runs agree with
the closed forms to within $3.1\times10^{-10}$; the $9$ negative controls are all rejected. The uniformised
law is a Poisson mixture of the $P^n$ \tierC. These computations therefore test the same formulas, but not the
Lean matrix powers, and \emph{they establish nothing}.
\end{numbox}

\subsection{Interpretation: does a resistance mutation take over the population?}

\begin{physbox}[: biological reading]
\tierC{} The Moran model is often read as a model of the following question: a bacterium or a cell becomes
resistant to a drug; does its lineage take over? If the mutation confers, in the presence of the
drug, a relative advantage $r>1$, the probability that a single mutant fixes is
\[
\rho_N(1)=\frac{1-r^{-1}}{1-r^{-N}}\xrightarrow[N\to\infty]{}1-\frac1r .
\]
The closed form is \tierA{} (Theorem~\ref{ch4:thm:moran}) and the limit in $N$ is an elementary computation \tierC.
Evaluated in floating point \tierX, it gives for $N=100$:
\begin{center}\small
\begin{tabular}{lccc}
\toprule
& $r=0.99$ & $r=1.01$ & $r=1.1$\\
\midrule
$\rho_{100}(1)$ & $0.0058$ & $0.0157$ & $0.0909$\\
neutral value $1/N$ & $0.01$ & $0.01$ & $0.01$\\
\bottomrule
\end{tabular}
\end{center}
Two qualitative lessons follow. First, even an advantageous mutant ($r=1.1$) is lost most of the time, here
in about $91\%$ of cases, swept away by drift. Second, a disadvantageous mutant ($r<1$) still has a small
chance of fixing.

\emph{Caution.} This is only a reading, and it remains at Tier~C. The model assumes a population of fixed size,
well mixed, a constant fitness, no new mutation, uniform death and generations that overlap step by step.
A real infection has a varying size, changing drug concentrations, a cost of resistance, spatial structure
and recurrent mutations. No data were used, and no Lean theorem certifies that a real population follows
this chain.
\end{physbox}

% ---------------------------------------------------------------------
\section{The Price equation and the haploid corollary}\label{ch4:sec:price}\sectionmark{The Price equation}

\subsection{An accounting identity}
A finite population consists of types $i\in s$. Type $i$ has a count $n_i$, a fitness $w_i$
(mean number of offspring per individual), a parental trait $z_i$, and its offspring have mean trait
$z'_i$. We write
\[
M=\sum n_i,\quad A=\sum n_iw_i,\quad B=\sum n_iz_i,\quad C=\sum n_iw_iz_i,\quad D=\sum n_iw_iz'_i .
\]
We then define $\bar w=A/M$ and $\bar z=B/M$. The mean of the next generation is $\bar z'=D/A$, since
type $i$ leaves $n_iw_i$ offspring: this is the \emph{modelling assumption}. Finally we write
$\mathrm{Cov}_n(w,z)=\frac1M\sum n_i(w_i-\bar w)(z_i-\bar z)$ and $\mathbb E_n[f]=\frac1M\sum n_if_i$.

\begin{theorem}[Price equation]\label{ch4:thm:price}
Over an arbitrary field, if $M\ne0$ and $A\ne0$, then
\[
\bar w\,(\bar z'-\bar z)=\mathrm{Cov}_n(w,z)+\mathbb E_n\bigl[w\,(z'-z)\bigr].
\]
\tierA{} (\lean{price_equation}, \texttt{formal\_cert/Programmes/PopGen/Price.lean}, ll.~95--105; real version
under positivity hypotheses \lean{price_equation_pos}, l.~170).
\end{theorem}

\begin{humanproof}
On the one hand, $\mathrm{Cov}_n(w,z)=\frac1M\bigl(C-\bar wB-\bar zA+M\bar w\bar z\bigr)=\frac{C-AB/M}{M}$; on the other
hand, $\mathbb E_n[w(z'-z)]=\frac{D-C}{M}$. The sum is $\frac DM-\frac{AB}{M^2}$. But
\[
\bar w(\bar z'-\bar z)=\frac AM\Bigl(\frac DA-\frac BM\Bigr)=\frac DM-\frac{AB}{M^2}. \qedhere
\]
\end{humanproof}

\begin{leanbox}[: the statement (\texttt{Price.lean}, ll.~95--98)]
\begin{lstlisting}[language=Lean,escapeinside={(*@}{@*)}]
theorem price_equation (s : Finset (*@$\iota$@*)) (n w z z' : (*@$\iota$@*) → K)
    (hN : ∑ i ∈ s, n i ≠ 0) (hW : ∑ i ∈ s, n i * w i ≠ 0) :
    wmean s n w * (zbarNext s n w z' - wmean s n z)
      = wcov s n w z + wmean s n (fun i => w i * (z' i - z i)) := by
\end{lstlisting}
The Lean proof follows the human proof. It expands the two numerators into the five raw sums
(\lean{cov_num_expand}, \lean{trans_num_expand}), then concludes with a scalar lemma in which $M,A,B,C,D$ are
atoms (\lean{price_scalar}: \lean{field_simp} then \lean{ring}). The statement holds over an arbitrary field \lean{K},
for example $\mathbb Q$ or $\R$, and for every finite set of types.
\end{leanbox}

\subsection{The haploid corollary}
\begin{corollary}[Haploid selection increment]\label{ch4:cor:haploide}
Take as trait the fitness itself, $z=w$, and assume there is \emph{no transmission
bias}, $z'=w$: clonal types, constant fitnesses, no mutation. Then
\[
\bar w\,(\bar w'-\bar w)=\mathrm{Var}_n(w),\qquad\text{hence}\qquad \Delta\bar w=\frac{\mathrm{Var}_n(w)}{\bar w}\ge0
\quad(\text{if }n_i>0,\ w_i>0).
\]
\tierA{} (\lean{fisher_haploid}, \lean{fisher_haploid_div}, \lean{fisher_haploid_nondecreasing},
\lean{fisher_haploid_pos}, \texttt{Price.lean}, ll.~125--185; \lean{fisher_haploid_nondecreasing} is re-exported under the name \lean{AURELIA}).
\end{corollary}

\begin{humanproof}
With $z=z'=w$, the transmission term $\mathbb E_n[w(w-w)]$ vanishes identically, and
$\mathrm{Cov}_n(w,w)=\mathrm{Var}_n(w)$. We divide by $\bar w\ne0$. The sign comes from $\mathrm{Var}_n(w)\ge0$ and
$\bar w>0$.
\end{humanproof}

This is equation~(2.27) of LaGatta, with a vanishing environmental term \tierL{}~\cite{ch4:LaG}. It is also
the discrete analogue of Baez's remark: $d\bar f/dt=\mathrm{Var}(f)$ when the fitnesses are
constant \tierL{}~\cite{ch4:Bae}.

For two alleles with frequencies $p$, $1-p$ and fitnesses $w_1$, $w_2$, one finds
$\mathrm{Var}_n(w)=p(1-p)(w_1-w_2)^2$, whence $\Delta\bar w=p(1-p)(w_1-w_2)^2/\bar w$. Module P7 derives from it, in
the case \lean{Fin 2}, two statements whose text is identical to two lemmas of the Genetic repository
(\texttt{Price.lean}, ll.~224 and 235)
\tierA. The connection between the two repositories remains textual \tierC.

\begin{remark}[Controls \tierA]
\emph{Positive control.} With three types, $n=(1,2,3)$, $w=(1,2,4)$, $z=(0,1,3)$ and $z'=(1,1,2)$, both sides
of the Price equation equal $-\frac{13}{36}$ (\lean{ctl_pos_three_types}, l.~249).
\emph{Transmission bias.} Take $n=(1,1)$ and $w=z=(1,2)$, but an offspring fitness
$z'=(1,1)$. Then $\Delta\bar w=-\frac12$, whereas $\mathrm{Var}/\bar w=\frac16$: the variance formula
\emph{fails}. The Price equation, however, still holds, with $-\frac34=-\frac34$ (\lean{ctl_neg_transmission_bias},
l.~263).
\emph{The factor $1/\bar w$ matters.} Without bias, $\Delta\bar w=\frac16$ but $\mathrm{Var}(w)=\frac14$
(\lean{ctl_neg_missing_wbar}, l.~280).
\end{remark}

\begin{warnbox}[: a negative result, announced as such]
The preregistration of programme P7 included an item \texttt{fisher\_general}: Fisher's fundamental theorem
in Price's (1972) ``partial change'' reading, with a prior probability of $0.35$. This
item was \textbf{not reached}. Only the haploid case with constant fitnesses is proved. Nor is it
the multi-allelic diploid case, where monotonicity is an inequality and not this equality. Price (1970,
1972) was not read; it is cited after LaGatta \cite{ch4:LaG} (p.~69).
\end{warnbox}

% ---------------------------------------------------------------------
\section{Evolutionary graph theory: the isothermal theorem}\label{ch4:sec:graphes}\sectionmark{The isothermal theorem}

\subsection{A structured population}
In a tissue, a dividing cell does not replace just any other cell: it replaces a
\emph{neighbour}. Evolutionary graph theory describes this structure by a weighted directed graph. The
individuals are the vertices $V=\{1,\dots,n\}$. The \emph{replacement matrix} $W$ has entries $\ge0$, and is
\emph{row-stochastic}: $\sum_jW_{ij}=1$. The entry $W_{ij}$ is the probability that the offspring of
$i$ replaces $j$; loops $W_{ii}>0$ are allowed.

A \emph{state} is the set $S\subseteq V$ of mutant vertices. Mutants have fitness $r>0$ and
residents fitness~$1$; we write $f_S(i)=r$ if $i\in S$ and $f_S(i)=1$ otherwise. Fix
selection weights $\mu_i>0$, and set $F(S)=\sum_k\mu_kf_S(k)$. One step of the \emph{birth--death} update
proceeds as follows.
\begin{itemize}[itemsep=1pt]
\item Vertex $i$ is chosen to reproduce with probability $\mu_if_S(i)/F(S)$.
\item Its offspring replaces $j$ with probability $W_{ij}$.
\item The new state is $S\cup\{j\}$ if $i\in S$, and $S\setminus\{j\}$ otherwise.
\end{itemize}
The case $\mu\equiv1$ is the process of Lieberman, Hauert and Nowak \cite{ch4:LHN}. The chain has $2^n$ states; $\emptyset$
and $V$ are absorbing. We write $P^{+}(S)$ and $P^{-}(S)$ for the probabilities that the \emph{number} of mutants $|S|$
increases, resp.\ decreases, by one.

\begin{definition}
The matrix $W$ is \emph{isothermal} if it is, in addition, column-stochastic: $\sum_iW_{ij}=1$. The
``temperature'' $\sum_iW_{ij}$ of each vertex then equals~$1$. It is \emph{strongly connected} if, for all
$a,b$, there is a path $a=x_0,\dots,x_m=b$, $m\ge1$, with $W_{x_kx_{k+1}}>0$.
\end{definition}

\begin{leanbox}[: the model (\texttt{formal\_cert/Programmes/EvoGraph/BirthDeath.lean}, ll.~44--61, 73--74, 90--91)]
\begin{lstlisting}[language=Lean]
noncomputable def fit (r : ℝ) (S : Finset (Fin n)) (i : Fin n) : ℝ := if i ∈ S then r else 1
noncomputable def ind (S : Finset (Fin n)) (i : Fin n) : ℝ := if i ∈ S then 1 else 0
def upd (S : Finset (Fin n)) (i j : Fin n) : Finset (Fin n) := if i ∈ S then insert j S else S.erase j
noncomputable def totFit (μ : Fin n → ℝ) (r : ℝ) (S : Finset (Fin n)) : ℝ := ∑ k, μ k * fit r S k
noncomputable def coef (μ : Fin n → ℝ) (W : Fin n → Fin n → ℝ) (r : ℝ) (S : Finset (Fin n))
    (i j : Fin n) : ℝ :=
  μ i * fit r S i / totFit μ r S * W i j
noncomputable def bdK (μ : Fin n → ℝ) (W : Fin n → Fin n → ℝ) (r : ℝ) (S T : Finset (Fin n)) : ℝ :=
  ∑ i, ∑ j, coef μ W r S i j * (if upd S i j = T then 1 else 0)
noncomputable def upP (μ : Fin n → ℝ) (W : Fin n → Fin n → ℝ) (r : ℝ) (S : Finset (Fin n)) : ℝ :=
  ∑ T, bdK μ W r S T * (if T.card = S.card + 1 then 1 else 0)
def StronglyConnected (W : Fin n → Fin n → ℝ) : Prop :=
  ∀ a b, Relation.TransGen (fun x y => 0 < W x y) a b
\end{lstlisting}
The independent auditor judged that \lean{upP}, \lean{downP} and \lean{bdK} correctly model the
Lieberman--Hauert--Nowak process. This is an AI judgement, not a human audit \tierC.
\end{leanbox}

\subsection{Forward bias, configuration by configuration}
For $S\subseteq V$, we define the \emph{cuts}
\[
\mathrm{out}(S)=\sum_{i\in S}\sum_{j\notin S}\mu_iW_{ij},\qquad \mathrm{in}(S)=\sum_{i\notin S}\sum_{j\in S}\mu_iW_{ij}.
\]

\begin{lemma}[Moves are cuts]\label{ch4:lem:cuts}
We have $P^{+}(S)=r\,\mathrm{out}(S)/F(S)$ and $P^{-}(S)=\mathrm{in}(S)/F(S)$. The chain is stochastic, and
$\emptyset$, $V$ are absorbing. \tierA{} (\lean{up_eq}, \lean{down_eq}, \lean{bd_absorbing},
\texttt{BirthDeath.lean}, ll.~148--250).
\end{lemma}

\begin{humanproof}
The coefficients $c(S,i,j)=\mu_if_S(i)W_{ij}/F(S)$ are $\ge0$, and their sum is $1$, because $\sum_jW_{ij}=1$ and
$\sum_i\mu_if_S(i)=F(S)$. There are four cases according to whether $i$ and $j$ are in $S$ or not.
\begin{itemize}[itemsep=1pt]
\item $i\in S$ and $j\notin S$: a mutant invades a resident, and $|S|$ increases. The weight is $\mu_irW_{ij}/F$.
\item $i\notin S$ and $j\in S$: $|S|$ decreases. The weight is $\mu_iW_{ij}/F$.
\item In the other two cases, the state does not change.
\end{itemize}
Summing, we obtain the two formulas. Finally, $\mathrm{upd}(\emptyset,i,j)=\emptyset$ and $\mathrm{upd}(V,i,j)=V$.
\end{humanproof}

\begin{proposition}[Cut balance and forward bias]\label{ch4:prop:bias}
Assume $\mu$ is stationary: $\sum_i\mu_iW_{ij}=\mu_j$ for every $j$. Then $\mathrm{out}(S)=\mathrm{in}(S)$ for
\emph{every} configuration $S$, and therefore
\[
P^{+}(S)=r\,P^{-}(S)\qquad\text{for every }S\subseteq V .
\]
In particular, this is the case if $W$ is isothermal and $\mu\equiv1$. \tierA{} (\lean{flow_balance},
\lean{cut_balance}, \lean{forward_bias}, \texttt{BirthDeath.lean}, ll.~176--256).
\end{proposition}

\begin{humanproof}
We compute the same quantity in two ways. On the one hand, $\sum_{i\in S}\sum_{j\in V}\mu_iW_{ij}=\mu(S)$, since
$\sum_jW_{ij}=1$. On the other hand, $\sum_{j\in S}\sum_{i\in V}\mu_iW_{ij}=\mu(S)$, by stationarity. Both double
sums contain the common part $\sum_{i\in S}\sum_{j\in S}\mu_iW_{ij}$. Once it is removed, what remains is
$\mathrm{out}(S)=\mathrm{in}(S)$. Lemma~\ref{ch4:lem:cuts} then gives $P^{+}=r\,\mathrm{out}/F=r\,\mathrm{in}/F=rP^{-}$.
If $W$ is isothermal, $\sum_i1\cdot W_{ij}=1=\mu_j$: this is the stationarity of $\mu\equiv1$.
\end{humanproof}

\begin{leanbox}[: forward bias (\texttt{BirthDeath.lean}, ll.~251--253)]
\begin{lstlisting}[language=Lean]
theorem forward_bias {μ : Fin n → ℝ} {W : Fin n → Fin n → ℝ} (r : ℝ) (hrow : ∀ i, ∑ j, W i j = 1)
    (hstat : ∀ j, ∑ i, μ i * W i j = μ j) (S : Finset (Fin n)) :
    upP μ W r S = r * downP μ W r S := by
\end{lstlisting}
The forward bias of Pattni et al. \tierL{}~\cite{ch4:PBRS} (Lemma~1, p.~11) and the ratio $p^-/p^+=1/r$ of
Theorem~3.2 of Keller--Uğurlu \tierL{}~\cite{ch4:KU} say the same thing. Note that the identity is written
without division: it remains true even if $P^{-}(S)=0$.
\end{leanbox}

\begin{proposition}[A harmonic function on configurations]\label{ch4:prop:harmgraph}
Let $\rho_n(k)=S_k/S_n$ be the Moran fixation function of Theorem~\ref{ch4:thm:moran}. If $\mu$ is
stationary, the function $h(S)=\rho_n(|S|)$ is harmonic for the configuration chain at every
$S\ne\emptyset$. If moreover an edge $i\in S\to j\notin S$ satisfies $\mu_iW_{ij}>0$, the one-step variance of $h$ at
$S$ is strictly positive. \tierA{} (\lean{moranFix_card_harmonic}, \lean{moranFix_card_var_pos},
\texttt{BirthDeath.lean}, ll.~285--328).
\end{proposition}

\begin{humanproof}
Set $k=|S|\ge1$. By Step~1 of Theorem~\ref{ch4:thm:moran}, $h$ goes up by $\Delta^{+}=r^{-k}/S_n$ on a
$+1$ step, and goes down by $\Delta^{-}=r^{-(k-1)}/S_n=r\,\Delta^{+}$ on a $-1$ step. Otherwise, $h$ does not change. Hence
\[
\mathbb E\bigl[h(S')-h(S)\bigr]=P^{+}(S)\Delta^{+}-P^{-}(S)\Delta^{-}=\frac{r^{-k}}{S_n}\bigl(P^{+}(S)-r\,P^{-}(S)\bigr)=0
\]
by Proposition~\ref{ch4:prop:bias}. For the variance, the move $i\to j$ has probability $>0$ and changes
$h$ by $\Delta^{+}\ne0$.
\end{humanproof}

This computation uses \emph{only} the forward bias. It never asserts that the number of mutants is a Markov
chain. We shall see in §\ref{ch4:ssec:cycle} that this assertion would be false.

\begin{lemma}[Connectivity forces cuts]\label{ch4:lem:conn}
Assume $W$ is strongly connected and $\mu>0$. Then every nonempty subset $S$ different from $V$ has an outgoing
edge $i\in S\to j\notin S$ with $\mu_iW_{ij}>0$. \tierA{} (\lean{cutPositive_of_strongly_connected},
\texttt{BirthDeath.lean}, l.~332).
\end{lemma}

\begin{humanproof}
Take $a\in S$, $b\notin S$, and a path $a=x_0,\dots,x_m=b$ of weight $>0$. Let $\ell$ be the smallest index such
that $x_{\ell+1}\notin S$. It exists because $x_m\notin S$, and $x_\ell\in S$, because $x_0\in S$. The edge
$x_\ell\to x_{\ell+1}$ leaves $S$. In Lean, the proof proceeds by induction on the \lean{TransGen} derivation.
\end{humanproof}

\subsection{The theorem}
\begin{theorem}[Isothermal theorem, birth--death update]\label{ch4:thm:iso}
Let $n\ge1$, $r>0$, and $W$ isothermal and strongly connected, with $\mu\equiv1$. For \emph{every} initial
configuration $S$,
\[
P^t(S,V)\xrightarrow[t\to\infty]{}\rho_n(|S|)=\frac{1-r^{-|S|}}{1-r^{-n}}\quad(r\ne1),\qquad\frac{|S|}{n}\quad(r=1).
\]
In particular, a single mutant placed at \emph{any} vertex $v$ fixes with limiting probability
$\frac{1-r^{-1}}{1-r^{-n}}$, exactly as in the well-mixed population. More generally, the same limit
holds as soon as $\mu W=\mu$, $\mu>0$, and every nonempty proper subset has an outgoing edge.
\tierA{} (\lean{isothermal_fixation}, \lean{isothermal_fixation_single}, \lean{isothermal_fixation_neutral},
\lean{isothermal_fixation_uniform}, \lean{stationary_fixation}, \texttt{BirthDeath.lean}, ll.~356--416; \lean{isothermal_fixation_single} is re-exported under the name \lean{ILENE}).
\end{theorem}

\begin{humanproof}
Proposition~\ref{ch4:prop:harmgraph} provides a function harmonic off $\{\emptyset,V\}$, equal to $0$ at
$\emptyset$ and $1$ at $V$. Lemma~\ref{ch4:lem:conn} and Proposition~\ref{ch4:prop:harmgraph} ensure that its
variance is $>0$ off the two absorbing states. The general absorption theorem of §\ref{ch4:sec:wf}
applies, on an arbitrary finite state space. We transport it there by numbering the $2^n$ configurations from
$0$ to $2^n-1$, with $\emptyset\mapsto0$ and $V\mapsto2^n-1$.
\end{humanproof}

\begin{leanbox}[: the isothermal theorem (\texttt{BirthDeath.lean}, ll.~376--378)]
\begin{lstlisting}[language=Lean]
theorem isothermal_fixation {W : Fin n → Fin n → ℝ} {r : ℝ} (hn : 0 < n) (hr : 0 < r)
    (hW : Isothermal W) (hsc : StronglyConnected W) (S : Finset (Fin n)) :
    Tendsto (fun t => sPow (bdK (fun _ => 1) W r) t S univ) atTop (𝓝 (moranFix n r S.card)) := by
\end{lstlisting}
\emph{Reuse.} The module \texttt{Transport.lean} reproves nothing. The theorem
\lean{two_absorbing_limit} (l.~81) conjugates the configuration chain by a bijection with
\lean{Fin (M + 1)}, then applies the theorem \lean{ProgPopGen.fixation_limit} of P7 \emph{as is}. The files
\texttt{Absorption.lean} and \texttt{MoranModel.lean} are reused without modification.
\emph{What is not proved.} The converse is not formalised: is a graph that gives the Moran value for
every $r$ necessarily isothermal? Proposition~2 of Pattni et al.\ is an equivalence; only the ``sufficient'' direction
is formalised here \tierL{}~\cite{ch4:PBRS}. The death--birth update is not formalised either.
\end{leanbox}

\subsection{Why the number of mutants does not ``lump''}\label{ch4:ssec:cycle}
\begin{center}
\setlength{\unitlength}{1cm}
\begin{picture}(9.5,3.0)
\put(0.5,2.5){\circle{0.5}}\put(0.42,2.42){\small 0}
\put(2.5,2.5){\circle{0.5}}\put(2.42,2.42){\small 1}
\put(2.5,0.5){\circle{0.5}}\put(2.42,0.42){\small 2}
\put(0.5,0.5){\circle{0.5}}\put(0.42,0.42){\small 3}
\put(0.75,2.5){\line(1,0){1.5}}\put(2.5,0.75){\line(0,1){1.5}}
\put(0.75,0.5){\line(1,0){1.5}}\put(0.5,0.75){\line(0,1){1.5}}
\put(3.3,1.4){\parbox{6cm}{\small The 4-cycle: each vertex sends its offspring to one of its two neighbours,
with probability $\frac12$. $W$ is doubly stochastic, hence isothermal.}}
\end{picture}
\end{center}

\begin{proposition}[The 4-cycle \tierA]\label{ch4:prop:cycle}
On the 4-cycle, with $r=2$ and $\mu\equiv1$, the configurations $\{0,1\}$ and $\{0,2\}$ both have two mutants,
but
\[
P^{+}(\{0,1\})=\tfrac13,\qquad P^{+}(\{0,2\})=\tfrac23,
\]
and the bias is $P^{+}=2P^{-}$ in both cases (\lean{cycle4_not_lumpable}, \texttt{Examples.lean},
ll.~239--253).
\end{proposition}

\begin{humanproof}
For $S=\{0,1\}$, we have $F=2+2+1+1=6$. The outgoing edges are $0\to3$ and $1\to2$, each of weight $\frac12$,
so $\mathrm{out}=1$, whence $P^{+}=2\cdot1/6=\frac13$. Likewise $\mathrm{in}=1$ and $P^{-}=\frac16$. For
$S=\{0,2\}$, the four edges leaving $0$ and $2$ all exit $S$, so $\mathrm{out}=\mathrm{in}=2$,
$P^{+}=\frac46=\frac23$ and $P^{-}=\frac13$.
\end{humanproof}

\begin{leanbox}[: the counterexample (\texttt{Examples.lean}, ll.~230--231 and 239--243)]
\begin{lstlisting}[language=Lean]
noncomputable def Wcyc : Fin 4 → Fin 4 → ℝ :=
  ![![0, 1 / 2, 0, 1 / 2], ![1 / 2, 0, 1 / 2, 0], ![0, 1 / 2, 0, 1 / 2], ![1 / 2, 0, 1 / 2, 0]]
theorem cycle4_not_lumpable :
    ({0, 1} : Finset (Fin 4)).card = ({0, 2} : Finset (Fin 4)).card ∧
    upP (fun _ => 1) Wcyc 2 {0, 1} = 1 / 3 ∧ upP (fun _ => 1) Wcyc 2 {0, 2} = 2 / 3 ∧
    upP (fun _ => 1) Wcyc 2 {0, 1} = 2 * downP (fun _ => 1) Wcyc 2 {0, 1} ∧
    upP (fun _ => 1) Wcyc 2 {0, 2} = 2 * downP (fun _ => 1) Wcyc 2 {0, 2} := by
\end{lstlisting}
\end{leanbox}

The transition probabilities of the count $|S_t|$ are therefore \emph{not} functions of $|S_t|$. The number of
mutants is not a birth--death chain with fixed rates, in the sense of the strong lumpability
of Kemeny and Snell \cite{ch4:KS}. Two clarifications are in order \tierC. First, we do not claim that the
number of mutants is non-Markovian for \emph{every} initial distribution: only strong lumpability is refuted.
Second, the \emph{jump chain}, the one in which the steps where nothing happens are removed, does lump. On the
4-cycle, given that a move takes place, both configurations go up with the same probability
$P^{+}/(P^{+}+P^{-})=\frac23$, whereas their probabilities of staying put are $\frac12$ and $0$ (exact rational
computation \tierX). This is the argument of Pattni et al. \tierL{}~\cite{ch4:PBRS}. The Lean proof never lumps
anything: it works configuration by configuration (Proposition~\ref{ch4:prop:harmgraph}).

\subsection{The star graph: amplifying or suppressing selection}
\begin{center}
\setlength{\unitlength}{1cm}
\begin{picture}(7,1.6)
\put(0.5,0.8){\circle{0.5}}\put(0.42,0.72){\small 1}
\put(3.5,0.8){\circle{0.5}}\put(3.42,0.72){\small 0}
\put(6.5,0.8){\circle{0.5}}\put(6.42,0.72){\small 2}
\put(3.25,0.95){\vector(-1,0){2.5}}\put(1.8,1.1){\small $\frac12$}
\put(0.75,0.65){\vector(1,0){2.5}}\put(1.9,0.25){\small $1$}
\put(3.75,0.95){\vector(1,0){2.5}}\put(4.9,1.1){\small $\frac12$}
\put(6.25,0.65){\vector(-1,0){2.5}}\put(5.0,0.25){\small $1$}
\end{picture}
\[
W_\star=\begin{pmatrix}0&\frac12&\frac12\\[1pt]1&0&0\\1&0&0\end{pmatrix}\qquad
(\text{centre }0,\ \text{leaves }1,2).
\]
\end{center}

The star is row-stochastic, but its column~$0$ sums to $2$: the centre has ``temperature'' $2$, and
the graph is \emph{not} isothermal (\lean{star_not_isothermal}) \tierA. Starting from a mutant at the centre, we have
$P^{+}(\{0\})=P^{-}(\{0\})=\frac12$. The ratio is $1$, not $r=2$ (\lean{star_bias}) \tierA.

\begin{proposition}[Exact values on the star \tierA]\label{ch4:prop:star}
For $n=3$, $r=2$ and $\mu\equiv1$, the limiting fixation probabilities are
\[
\tfrac5{12}\ \text{(mutant at the centre)},\qquad \tfrac23\ \text{(mutant on a leaf)},\qquad
\tfrac7{12}\ \text{(mutant placed uniformly)},
\]
all different from the Moran value $\rho_3(1)=\frac47$ (\lean{star_fixation_centre},
\lean{star_fixation_leaf}, \lean{star_fixation_uniform}, \lean{star_ne_moran}, \texttt{Examples.lean},
ll.~171--196). With the stationary weights $\mu=(\frac12,\frac14,\frac14)$, one recovers $\frac47$ starting from the
centre (\lean{starStationary_fixation}, l.~219).
\end{proposition}

The proof by hand is the subject of Exercise~\ref{ch4:ex:star}. In Lean, the table of the $2^3$ values
$h=(0,\frac5{12},\frac23,\frac23,\frac56,\frac56,\frac{14}{15},1)$ is given explicitly (\lean{hStar}, l.~139).
One checks that it is harmonic (\lean{star_harmonic}) and has positive variance, then feeds it into the
\emph{same} transported absorption machinery. An independent script in exact rational arithmetic gives the same table
\tierX.

\begin{physbox}[: amplifiers and suppressors]
\tierC{} Compare with the Moran value $\frac47\approx0.571$. For a mutant placed uniformly at random, the star
\emph{amplifies} selection, since $\frac7{12}\approx0.583>\frac47$. For a mutant born at the centre, it
\emph{suppresses} it, with $\frac5{12}\approx0.417$. These conclusions must be read with care. They are only
\emph{three numbers}, for $n=3$ and $r=2$. They do not show that the star amplifies selection for every $r$ or
for every $n$. The preregistration, incidentally, called the star an ``amplifier'' without qualification. For $n=4$,
an unverified script gives in exact rationals $\frac{392}{1265}$ (centre), $\frac{168}{253}$ (leaf) and a
mean $\frac{728}{1265}\ne\frac8{15}$. These values are \tierX{} and are not checked by the
kernel. We do not quote the classical formula for the star, because the source (Lieberman--Hauert--Nowak) was not
read.
\end{physbox}

\subsection{Connectivity is necessary}
\begin{proposition}[Identity graph \tierA]\label{ch4:prop:identite}
On $2$ vertices, the matrix $W=I$ is doubly stochastic, hence isothermal, but not strongly connected. With
$r=2$, we have $P^t(\{0\},V)=0$ for every $t$, which therefore does not tend to $\rho_2(1)=\frac23$
(\lean{identity_no_fixation}, \texttt{Examples.lean}, ll.~289--292).
\end{proposition}
Each individual replaces only itself, and nothing ever moves. The strong connectivity hypothesis of
Theorem~\ref{ch4:thm:iso} therefore cannot be dropped. This is a \emph{counterexample}, not a proved
converse \tierC.

\subsection{A preregistration error, told honestly}
\begin{warnbox}[: what was announced, and what was proved]
\tierC{} On 26 October 2026, item N4 of the preregistration announced: prove that if the graph is isothermal,
``\emph{the number of mutants is a birth-death chain with the Moran ratio $r$}'', and deduce from it the fixation
$(1-r^{-1})/(1-r^{-N})$ by reusing P7. Here is what happened.
\begin{itemize}[itemsep=1pt]
\item The \emph{conclusion} is proved, and even more. It holds for every initial configuration, for a
mutant at any vertex, for a uniformly placed mutant, and for $r=1$ (Theorem~\ref{ch4:thm:iso}).
\item The italicised \emph{intermediate clause} is \textbf{false} for the chain that includes the steps where nothing
happens. It was therefore not formalised. The 4-cycle counterexample (Proposition~\ref{ch4:prop:cycle}) is
checked by the kernel.
\item What is true and proved is the forward bias \emph{configuration by configuration}:
$P^{+}(S)=rP^{-}(S)$ (Proposition~\ref{ch4:prop:bias}).
\end{itemize}
This is not a discovery. Keller and Uğurlu observe that the Lieberman--Hauert--Nowak argument projects the
chain onto the number of mutants, and that ``the projection does not need to be Markovian'' \tierL{}~\cite{ch4:KU}
(abstract, p.~4). The preregistration error thus reproduced a known gap in a published proof. Keller and
Uğurlu also write that the theorem ``remains unproven till today''. We do not repeat this sentence, because
Lemma~1 of Pattni et al., which is earlier, gives a valid proof of it via the jump chain \tierC. The lesson is one of
method: state the goal, let the kernel decide, then publish the discrepancy.
\end{warnbox}

\begin{physbox}[: biological scope]
\tierC{} In this model, a balanced (doubly stochastic) and strongly connected structure does not change the
fixation probability of a mutant of fitness $r$, wherever it appears. A structure that amplifies or
suppresses selection must therefore break this balance; the star is an example of this, for a single choice of the
parameters. Tissues, where cells replace their neighbours, are the usual motivation. We make,
however, no claim specific to any tissue, and use no data. The model omits mutation,
other update rules, changing structures and frequency dependence.
\end{physbox}

% ---------------------------------------------------------------------
\section{Summary: what is verified, what is not}\label{ch4:sec:bilan}\sectionmark{Summary}

\begin{center}\small
\begin{tabular}{p{0.44\linewidth}p{0.33\linewidth}c}
\toprule
statement & Lean declaration (module) & tier\\
\midrule
Hardy--Weinberg stationary & \lean{hwStep_stationary} (H12) & A\\
increment identity, monotonicity of $\bar w$ & \lean{wbar_increment} (H12) & A\\
$i/N$ harmonic and unique (maximum principle) & \lean{harmonic_unique} (H11) & A\\
$P^n(i,N)\to h(i)$, general absorbing chain & \lean{fixation_limit} (P7) & A\\
Wright--Fisher: $P^n(i,N)\to i/N$ & \lean{wright_fisher_fixation_limit} (P7) & A\\
geometric decay of $I_n$ & --- (Lemma~\ref{ch4:lem:geom}) & C\\
Moran: $P^n(i,N)\to\frac{1-r^{-i}}{1-r^{-N}}$ & \lean{moran_fixation_limit_closed} (P7) & A\\
Price equation; $\Delta\bar w=\mathrm{Var}/\bar w$ (haploid) & \lean{price_equation}, \lean{fisher_haploid_div} (P7) & A\\
Fisher's theorem, ``partial change'' reading & not reached & ---\\
forward bias $P^{+}=rP^{-}$ & \lean{forward_bias} (N4) & A\\
isothermal theorem (sufficient direction) & \lean{isothermal_fixation} (N4) & A\\
4-cycle: $\frac13\ne\frac23$ & \lean{cycle4_not_lumpable} (N4) & A\\
star $n=3$: $\frac5{12}$, $\frac23$, $\frac7{12}$; Moran $\frac47$ & \lean{star_fixation_uniform} (N4) & A\\
star $n=4$ & exact script & X\\
``fixation probability'' $=\lim P^n(i,N)$ & no path space & C\\
biological reading (resistance, tissues) & --- & C\\
\bottomrule
\end{tabular}
\end{center}

\paragraph{Kernel replay performed for this chapter.}
On 2026-10-26, the writing agent ran \texttt{leanchecker} (Lean~4.34.1, pinned toolchain, two
cores) on the ten modules cited. They are \lean{Autoresearch.PopGenFisher}, \lean{Autoresearch.PopGenWrightFisher},
\lean{Programmes.PopGen.Absorption}, \lean{WrightFisherFixation}, \lean{MoranModel}, \lean{MoranFixation},
\lean{Price}, together with \lean{Programmes.EvoGraph.Transport}, \lean{BirthDeath} and \lean{Examples}. All ten
runs returned code~$0$, in $55$ to $131$~s each, with no output at all. The five modules previously without replay
(\lean{PopGenFisher}, \lean{PopGenWrightFisher}, \lean{WrightFisherFixation}, \lean{MoranFixation}, \lean{Price}) therefore now satisfy the
definition of Tier~A in this book. The \texttt{grep} search for \lean{skipKernelTC}, \lean{addDecl},
\lean{native_decide}, \lean{implemented_by}, \lean{sorry}, \lean{admit} and \lean{axiom} in these ten sources
finds only two mentions in documentation comments (\texttt{MoranModel.lean}, l.~237;
\texttt{Examples.lean}, l.~6). Three limitations should be noted.
\begin{itemize}[itemsep=1pt]
\item The replay was run on the producer's \texttt{.olean} files, not on a clean recompilation. These files
are all dated after their source.
\item The replay is not \texttt{--fresh}: Mathlib and the Lean core were loaded, not replayed.
\item The only negative control of this run is a nonexistent module, rejected with code~$1$. The control
by falsification under \lean{debug.skipKernelTC}, rejected by the same binary, is the one from the N4 audit.
\end{itemize}
Like the rest of the chapter, this verification is the work of an AI agent; no human carried it out. The
trace is \texttt{results/kernel\_replay\_2026-10-26/ch4\_popgen\_trace/leancheck\_ch4.json}. The replay of the whole repository, run afterwards
(\texttt{results/kernel\_replay\_2026-10-26/replay.jsonl}), also returned code~$0$ for \lean{PopGenFisher}, \lean{PopGenWrightFisher}
and \lean{PopGenImportBoth}.

\paragraph{What is not formalised.} The following points remain beyond the scope of this chapter.
\begin{itemize}[itemsep=1pt]
\item There is no measure on path space. A natural next step would be to identify \lean{kPow}
with the law constructed by \lean{ProbabilityTheory.Kernel.traj} (Ionescu--Tulcea, in Mathlib, after Marion
\tierL{}~\cite{ch4:Mar}).
\item Also not formalised: absorption times, Kimura's diffusion limit, mutation,
multi-allelic diploid selection, frequency dependence, the converse of the isothermal theorem and the
death--birth update.
\item No human has audited the adequacy of the statements. The ledgers of P7 and N4 flag
every Tier~A line with \texttt{LEDGER\_UNAUDITED\_TIER\_A}: there are $14$ of them for P7 and $11$
for N4. The Elenchus gate exits with code~$1$, and this is not hidden.
\end{itemize}

\paragraph{How to cite.} Programme P7 (fixation, Moran, Price) is published as doi:10.5281/zenodo.23653850, and
programme N4 (isothermal theorem) as doi:10.5281/zenodo.23660071 \cite{ch4:P7,ch4:N4}. Neither has been
peer reviewed.

% ---------------------------------------------------------------------
\section{Exercises}\label{ch4:sec:exercices}

\begin{exercise}[Hardy--Weinberg]\label{ch4:ex:hw}
A population has genotype frequencies $(x_{11},x_{12},x_{22})=(0.5,\,0.2,\,0.3)$.
\begin{enumerate}[label=\alph*),itemsep=0pt]
\item Compute $p$, then the genotype frequencies after one generation of random mating.
\item Check that a second generation changes nothing further. Which Lean theorem expresses this fact?
\end{enumerate}
\emph{Answer.} $p=0.6$, then $(0.36,\,0.48,\,0.16)$. The theorem is \lean{hwStep_stationary}.
\end{exercise}

\begin{exercise}[Genic selection]
Take $w_{11}=(1+s)^2$, $w_{12}=1+s$ and $w_{22}=1$, with $s>-1$.
\begin{enumerate}[label=\alph*),itemsep=0pt]
\item Show that $\bar w(p)=(1+ps)^2$ and that $p'=\dfrac{p(1+s)}{1+ps}$.
\item Deduce that, if $s>0$ and $0<p<1$, the sequence of frequencies is strictly increasing and converges to
$1$.
\end{enumerate}
\emph{Hint.} For b), write $1-p'=\frac{1-p}{1+ps}$ and iterate the inequality $1-p_{t+1}\le\frac{1-p_t}{1+p_0s}$.
\end{exercise}

\begin{exercise}[The increment identity]\label{ch4:ex:increment}
\begin{enumerate}[label=\alph*),itemsep=0pt]
\item Prove the key identity $\bar w+a\,pq=p\,w_{11}+q\,w_{22}$, with $a=w_{11}-2w_{12}+w_{22}$.
\item Deduce~\eqref{ch4:eq:increment} without copying the proof given in the text.
\item For $w=(2,1,1)$ and $p=\frac12$, compute $\bar w(p)$ and $\bar w(p')$. Check the identity on this example.
\item For $w=(-1,1,0)$ and $p=\frac12$, show that $\bar w$ \emph{decreases}. Which factor
of~\eqref{ch4:eq:increment} changes sign?
\item \emph{Underdominance.} With $w=(1,\frac12,1)$, show that $p=\frac12$ is an equilibrium. Compute $p'$ and
$\bar w(p')$ for $p=\frac25$, and comment.
\end{enumerate}
\emph{Answer.} c) $\bar w=\frac54$ and $\bar w(p')=\frac{34}{25}$. d) $\bar w=\frac14$, $p'=0$ and $\bar w(p')=0$;
the factor $\bar w+pw_{11}+qw_{22}=-\frac14$ is negative. e) $p'=\frac7{19}$, and $\bar w$ goes from $\frac{19}{25}$
to $\frac{277}{361}$: the mean goes up, but $p$ moves away from the equilibrium.
\end{exercise}

\begin{exercise}[Wright--Fisher with two genes]\label{ch4:ex:wf2}
Take $N=2$.
\begin{enumerate}[label=\alph*),itemsep=0pt]
\item Write down the matrix $P$ ($3\times3$).
\item Show that $P^n(1,1)=2^{-n}$ and that $P^n(1,2)=\frac12\bigl(1-2^{-n}\bigr)$. Recover $P^2(1,2)=\frac38$
(\lean{ctl_pos_wf_two_steps}).
\item Compare the rate obtained with the bound of Lemma~\ref{ch4:lem:geom}.
\end{enumerate}
\emph{Answer.} c) Here $\delta=\frac14$, so the bound is $(\frac34)^n$, whereas the true decay is
$(\frac12)^n$.
\end{exercise}

\begin{exercise}[Wright--Fisher with three genes]
Take $N=3$.
\begin{enumerate}[label=\alph*),itemsep=0pt]
\item Compute rows $1$ and $2$ of $P$. Check that, from each interior state, the probability of staying
in the interior is $\frac23$.
\item Deduce that $I_n(1)=(\frac23)^n$ exactly, and that $\bigl|P^n(1,3)-\frac13\bigr|\le\frac23\,(\frac23)^n$.
\end{enumerate}
\emph{Answer.} a) Row $1$ is $(\frac8{27},\frac49,\frac29,\frac1{27})$ and row $2$ is
$(\frac1{27},\frac29,\frac49,\frac8{27})$.
\end{exercise}

\begin{exercise}[Moran, $N=3$, $r=2$]\label{ch4:ex:moran3}
\begin{enumerate}[label=\alph*),itemsep=0pt]
\item Compute $p_1,q_1,p_2,q_2$ and check that $p_i=2q_i$.
\item Solve the gambler's-ruin recurrence, and recover $h=(0,\frac47,\frac67,1)$.
\item Show that the ``inverted'' formula $(1-r^i)/(1-r^N)$ gives $\frac17$ at $i=1$, and that it does not satisfy
the first-step equation.
\end{enumerate}
\emph{Answer.} a) $p_1=\frac13$, $q_1=\frac16$, $p_2=\frac4{15}$, $q_2=\frac2{15}$.
\end{exercise}

\begin{exercise}[Large population]
\begin{enumerate}[label=\alph*),itemsep=0pt]
\item If $r>1$, show that $\rho_N(1)=\frac{1-r^{-1}}{1-r^{-N}}$ decreases to $1-\frac1r$ as $N\to\infty$.
\item If $0<r<1$, show that $\rho_N(1)\sim(r^{-1}-1)\,r^N$. Compare with the neutral value $1/N$.
\item For $r=1.1$, what proportion of advantageous mutants is lost in a very large population? Why
does this reading remain at Tier~C?
\end{enumerate}
\end{exercise}

\begin{exercise}[Price equation]
\begin{enumerate}[label=\alph*),itemsep=0pt]
\item Take $n=(1,2,3)$, $w=(1,2,4)$, $z=(0,1,3)$ and $z'=(1,1,2)$. Compute $\bar w$, $\bar z$, $\bar z'$,
$\mathrm{Cov}_n(w,z)$ and $\mathbb E_n[w\Delta z]$, and check that both sides equal $-\frac{13}{36}$.
\item For two types with frequencies $p$, $1-p$, show that $\mathrm{Var}_n(w)=p(1-p)(w_1-w_2)^2$.
\item Take $n=(1,1)$, $w=z=(1,2)$ and $z'=(1,1)$. Show that $\Delta\bar w\ne\mathrm{Var}_n(w)/\bar w$, whereas
the Price equation still holds. Which hypothesis of Corollary~\ref{ch4:cor:haploide} is missing?
\end{enumerate}
\emph{Answer.} a) $\bar w=\frac{17}6$, $\bar z=\frac{11}6$ and $\bar z'=\frac{29}{17}$, so
$\bar z'-\bar z=-\frac{13}{102}$; b) follows by a direct
computation; c) $\Delta\bar w=-\frac12$ and $\mathrm{Var}/\bar w=\frac16$; what is missing is the absence of transmission
bias.
\end{exercise}

\begin{exercise}[The complete graph is the Moran model]
Take $W_{ij}=\frac1n$ for all $i,j$ (loops included) and $\mu\equiv1$.
\begin{enumerate}[label=\alph*),itemsep=0pt]
\item Check that $W$ is isothermal and strongly connected.
\item For $|S|=k$, compute $\mathrm{out}(S)$, $\mathrm{in}(S)$ and $F(S)$. Deduce that $P^{+}(S)$ and $P^{-}(S)$
are exactly the rates $p_k$, $q_k$ of the Moran model (\lean{wellMixed_up}, \lean{wellMixed_down},
\texttt{Examples.lean}, ll.~62--81).
\item Deduce that, here, the configuration chain lumps exactly onto the Moran chain.
\end{enumerate}
\emph{Answer.} b) $\mathrm{out}=\mathrm{in}=k(n-k)/n$ and $F=rk+n-k$.
\end{exercise}

\begin{exercise}[The 4-cycle, by hand]
\begin{enumerate}[label=\alph*),itemsep=0pt]
\item Check that the matrix of the 4-cycle is doubly stochastic.
\item Recover $P^{\pm}(\{0,1\})=(\frac13,\frac16)$ and $P^{\pm}(\{0,2\})=(\frac23,\frac13)$.
\item Compute, for each of the two configurations, the probability of going up \emph{given that a move takes
place}, and comment.
\item What does the isothermal theorem give for a single mutant, when $r=2$?
\end{enumerate}
\emph{Answer.} c) $\frac23$ in both cases. d) $\rho_4(1)=\frac{1-1/2}{1-1/16}=\frac8{15}$.
\end{exercise}

\begin{exercise}[Fixation probability on the star, by hand]\label{ch4:ex:star}
Take the three-vertex star (centre $0$, leaves $1$ and $2$), with $W_{0j}=\frac12$, $W_{j0}=1$, $r=2$ and
$\mu\equiv1$. By symmetry, set $x_0=h(\{0\})$, $x_1=h(\{1\})=h(\{2\})$, $y=h(\{0,1\})=h(\{0,2\})$ and
$z=h(\{1,2\})$.
\begin{enumerate}[label=\alph*),itemsep=0pt]
\item From $\{0\}$: show that $F=4$, that the chain moves to $\{0,j\}$ with probability $\frac12$ (in total) and to
$\emptyset$ with probability $\frac12$. Deduce that $x_0=\frac y2$.
\item From $\{1\}$: show that the chain moves to $\{0,1\}$ with probability $\frac12$, to $\emptyset$ with
probability $\frac18$, and stays put with probability $\frac38$. Deduce that $x_1=\frac45y$.
\item From $\{0,1\}$: show that the chain moves to $V$ with probability $\frac15$ and to $\{1\}$ with probability
$\frac15$, and deduce that $2y=1+x_1$. From $\{1,2\}$: show that $z=\frac45+\frac{x_1}5$.
\item Solve the system, and recover $x_0=\frac5{12}$, $x_1=\frac23$, $y=\frac56$ and $z=\frac{14}{15}$, then the
value $\frac7{12}$ for a uniformly placed mutant. Compare with $\frac47$.
\item Redo a) with $\mu=(\frac12,\frac14,\frac14)$. Check that $\mu W_\star=\mu$, then that $x_0=\frac47$
(\lean{starStationary_fixation}).
\end{enumerate}
\emph{Answer.} d) One obtains $y=\frac12+\frac25y$, whence $y=\frac56$.
\end{exercise}

\begin{exercise}[Reading a Lean statement]
\begin{enumerate}[label=\alph*),itemsep=0pt]
\item Take $N=2$. Show that the absorbing kernels for which $j\mapsto j/2$ is harmonic are
exactly those of the form $P(1,0)=P(1,2)=a$, $P(1,1)=1-2a$, with $a\in[0,\frac12]$. Show that the hypothesis
\lean{hvar} of \lean{fixation_limit} is equivalent to $a>0$, and that in this case $P^n(1,2)=\frac12\bigl(1-(1-2a)^n\bigr)$.
What happens for $a=0$?
\item Why does a file containing \lean{sorry} compile with return code $0$? What are the three
checks that make up the Tier~A gate in this book?
\item Does the statement \lean{isothermal_fixation} talk about probabilities? Write in one sentence the \tierC{} bridge that
connects it to the ``fixation probability''.
\end{enumerate}
\emph{Hint.} a) The one-step variance at $1$ is $2a\cdot\frac14=\frac a2$. For $a=0$, one recovers the
identity kernel, that is, the negative control \lean{ctl_neg_identity_no_absorption}.
\end{exercise}

\begin{exercise}[($\dagger$) University exercise: Expected absorption time via the fundamental matrix]
Consider the neutral Moran chain ($r=1$) with $N=3$, so that $p_i=q_i=i(N-i)/N^2$.
\begin{enumerate}[label=\alph*),itemsep=0pt]
\item Write down the interior block $Q$ (states $1,2$) and the interior-to-boundary block $R$ (columns $0,3$) of $P$.
\item Compute the fundamental matrix $(I-Q)^{-1}$, and deduce the expected absorption time $t_i=\mathbb E_i[T]$,
$T=\min(T_0,T_3)$, from $i=1$ and $i=2$, using $t=(I-Q)^{-1}\mathbf 1$.
\item Compute $B=(I-Q)^{-1}R$ and check that its column for state $3$ is $(\frac13,\frac23)$, in agreement with
Theorem~\ref{ch4:thm:moran} for $r=1$.
\end{enumerate}
\emph{Answer.} a) $Q=\begin{pmatrix}\frac59&\frac29\\[1pt]\frac29&\frac59\end{pmatrix}$,
$R=\begin{pmatrix}\frac29&0\\[1pt]0&\frac29\end{pmatrix}$. b) $(I-Q)^{-1}=\begin{pmatrix}3&\frac32\\[1pt]\frac32&3\end{pmatrix}$,
so $t_1=t_2=\frac92$. c) $B=\begin{pmatrix}\frac23&\frac13\\[1pt]\frac13&\frac23\end{pmatrix}$.
\end{exercise}

\begin{exercise}[($\dagger$) University exercise: The Moran fixation formula via a martingale]\label{ch4:ex:moranmart}
Let $(X_n)$ be the Moran chain with $N\ge1$ and $r>0$, $r\ne1$, started at $X_0=i$, and let $T=\min(T_0,T_N)$.
\begin{enumerate}[label=\alph*),itemsep=0pt]
\item Using $p_k=r\,q_k$ (Lemma~\ref{ch4:lem:moranbase}), show that
$\mathbb E\bigl[r^{-X_{n+1}}\mid X_n=k\bigr]=r^{-k}$ for every $k$, so that $M_n=r^{-X_n}$ is a martingale.
\item Justify that $T<\infty$ almost surely (for instance from $I_n(i)\to0$ in the proof of
Theorem~\ref{ch4:thm:moran}), and that $M$ is bounded.
\item Apply the optional stopping theorem to obtain $\mathbb P_i(X_T=N)=\dfrac{1-r^{-i}}{1-r^{-N}}$.
\item Which martingale plays the same role for $r=1$?
\end{enumerate}
\emph{Answer.} a) $\mathbb E[r^{-X_{n+1}}\mid X_n=k]-r^{-k}=r^{-k}(r-1)\bigl(q_k-p_k/r\bigr)=0$.
c) $r^{-i}=\mathbb E_i[r^{-X_T}]=u\,r^{-N}+(1-u)$ with $u=\mathbb P_i(X_T=N)$. d) $X_n$ itself, which gives $i/N$.
\end{exercise}

\label{ch4:fin}
\fussy
